\head{第四卷勘误清单三：与第0章§§16--23已译部分有关的选译}
\source{EGA IV第四部分，勘误清单三，印刷第345--348页；上传PDF第344--347页}
\NewNote{本节分批收入与已译范围直接相关的作者勘误，保留其原分类和原编号。它不是整份勘误清单三的全译；未列出的条目不因此被删除或判定无关。下文的“页”均指作者所列原出版分册的印刷页码。}

\egasection{A）排印错误：七项选译}
\hypertarget{ErrIV:A16222}{}\noindent\textbf{（$0_{\mathrm{IV}}$，16.2.2.2）。}
第26页第13行，将 $M'\cap\mathfrak q^nM'$ 改为 $M'\cap\mathfrak q^nM$。

\medskip
\hypertarget{ErrIV:A1737}{}\noindent\textbf{（$0_{\mathrm{IV}}$，17.3.7）。}
第50页倒数第16行，将（17.1.10）改为（17.3.3）。

\medskip
\hypertarget{ErrIV:A1813}{}\noindent\textbf{（$0_{\mathrm{IV}}$，18.1.3）。}
在图（18.1.3.1）中，$p_2$ 应标在箭头 $G\to F$ 的旁边。

\medskip
\hypertarget{ErrIV:A1987}{}\noindent\textbf{（$0_{\mathrm{IV}}$，19.8.7）。}
第111页倒数第7、8、9行，将（21.5.5）改为（21.5.3）。

\medskip
\hypertarget{ErrIV:A1988}{}\noindent\textbf{（$0_{\mathrm{IV}}$，19.8.8）。}
第112页倒数第2行，将“$A$-模”改为“$B$-模”。

\medskip
\hypertarget{ErrIV:A2054}{}\noindent\textbf{（$0_{\mathrm{IV}}$，20.5.4）。}
第130页倒数第5行，将 $\Omega_{B'/A'}$ 改为 $\Omega^1_{B'/A'}$。


\par\medskip
\hypertarget{ErrIV:A2078}{}\textbf{（$0_{\mathrm{IV}}$，20.7.8）}
第149页倒数第9行，将（18.4.3）改为（18.5.4）；倒数第4行，将（18.4.1）改为（18.5.1）。
\NewNote{这是作者的正式引文勘误。第十版20.7.8正文已经采用两个修正后的出处，原引文在脚注J10-C01中保留。}

\egasection{B）正文修订：十项选译}
\hypertarget{ErrIV:5}{}
\noindent\textbf{（Err$_{\mathrm{IV}}$，5）。}
在（$0_{\mathrm{IV}}$，16.1.4）中，第23页倒数第8行，将“说 $A$ 是悬链的”改为：“说一个整环 $A$ 是悬链的”。

\medskip
\hypertarget{ErrIV:6}{}
\noindent\textbf{（Err$_{\mathrm{IV}}$，6）。}
在（$0_{\mathrm{IV}}$，16.1.9）中，可删去 $A$ 为诺特环的假设。事实上，对于每个 $u\in B_M$，$A[u]$ 是 $C$ 的一个子 $A$-代数，故 $u$ 整于 $A$（Bourbaki，\textit{交换代数}，第五章，§1，第1目，命题1的推论1），因而也整于 $A_M$；仍由（16.1.5）得出结论。

\medskip
\hypertarget{ErrIV:7}{}
\noindent\textbf{（Err$_{\mathrm{IV}}$，7）。}
在（$0_{\mathrm{IV}}$，16.5.12）中，应注意：关系（16.5.11.1）等价于将 $A$ 换成各商环 $A/\mathfrak r$ 后得到的相同关系，其中 $\mathfrak r$ 遍历 $A$ 的极小素理想。因此，在（16.5.12）的证明中可限于 $A$ 为整环的情形，并应用（16.1.4.2）。


\medskip
\hypertarget{ErrIV:8}{}
\noindent\textbf{（Err$_{\mathrm{IV}}$，8）。}
删去（$0_{\mathrm{IV}}$，19.3.12），因为其证明错误（集合 $W_\lambda,U_\lambda,V_\lambda$ 不是线性流形）。


\medskip
\hypertarget{ErrIV:9}{}
\noindent\textbf{（Err$_{\mathrm{IV}}$，9）。}
在（$0_{\mathrm{IV}}$，19.8.3）中，第110页倒数第6、8行，将“同态”改为“局部同态”；在（$0_{\mathrm{IV}}$，19.8.4）第111页第1行同样修改。


\medskip
\hypertarget{ErrIV:10}{}\noindent\textbf{（Err$_{\mathrm{IV}}$，10）。}
在（$0_{\mathrm{IV}}$，20.1.2）中，第117页倒数第2行，将 $a\in A$ 改为 $b\in B$，将 $aD$ 改为 $bD$；倒数第1行，将“$A$-模”改为“$B$-模”。

\medskip
\hypertarget{ErrIV:11}{}\noindent\textbf{（Err$_{\mathrm{IV}}$，11）。}
在（$0_{\mathrm{IV}}$，20.2.1）中，第119页第8行，将“$A$-模”改为“$B$-模”；第21--22行，将“关于 $u$ 的模的双同态”改为“$B$-模同态”。

\medskip
\hypertarget{ErrIV:12}{}\noindent\textbf{（Err$_{\mathrm{IV}}$，12）。}
在（$0_{\mathrm{IV}}$，20.4.6.1）中，将 $p_1-p_2$ 改为 $p_2-p_1$；第125页倒数第10行，将 $x\otimes1-1\otimes x$ 改为 $1\otimes x-x\otimes1$。

\medskip
\hypertarget{ErrIV:13}{}\noindent\textbf{（Err$_{\mathrm{IV}}$，13）。}
在（$0_{\mathrm{IV}}$，20.4.8）中，第126页第14行，将 $p_2=p_1-d_{B/A}$ 改为 $p_1=p_2-d_{B/A}$。

\medskip
\hypertarget{ErrIV:14}{}\noindent\textbf{（Err$_{\mathrm{IV}}$，14）。}
在（$0_{\mathrm{IV}}$，20.4.13）中，将第127页倒数第6--11行替换如下：

$(iv)$ 设 $A,B$ 为整环，$A\subseteq B$，$B$ 整于 $A$，且 $B$ 的分式域是 $A$ 的分式域的可分扩张。则 $B$-模 $\Omega^1_{B/A}$ 为挠模。事实上，对于 $x\in B$，令 $f$ 为 $x$ 对 $A$ 的分式域 $K$ 的最小多项式；存在非零 $a\in A$，使 $af(T)\in A[T]$。由 $x$ 对 $K$ 可分，$f'(x)\ne0$；又对 $af(x)=0$ 微分得 $af'(x)d_{B/A}x=0$。结合（20.4.7），即得所述结论。


\subsec{本版补入的§21相关排印勘误}
\noindent 以下两项出自《勘误与增补（清单三）》A部，第四卷第四部分原分册第345页（上传PDF第344页）。原勘误页码指被订正的原分册。

\medskip
\noindent\textbf{（$0_{\mathrm{IV}}$，21.5.3）。}
第166页倒数第1行，将法文“ramifiée”改为“ramifié”。\footnote{译者说明：这是法文性数配合的修正，中文“形式非分歧”无须相应换字。}

\medskip
\noindent\textbf{（$0_{\mathrm{IV}}$，21.7.1）。}
第170页第4行，将（20.6.15.1）改为（20.6.17.1）。


\subsec{本版补入的22.1--22.3相关正式勘误}
以下三项出自《勘误与增补（清单三）》A部原分册第345--346页及B部第15项（第348页）；以下被订正页码均指第四卷第一部分的原分册印刷页码。

\medskip
\noindent\textbf{（$0_{\mathrm{IV}}$，22.1.1）。}
第183页倒数第12行，将（20.6.20）改为（20.6.10）。

\medskip
\noindent\textbf{（$0_{\mathrm{IV}}$，22.3.3）。}
第192页倒数第13行，将 $\widehat B_{\mathfrak p}$ 改为 $\widehat B_{\mathfrak q}$。

\medskip
\hypertarget{ErrIV:15}{}\noindent\textbf{（Err$_{\mathrm{IV}}$，15）。}
在（$0_{\mathrm{IV}}$，22.2.7）中，第189页第5行，将“形式光滑”改为“对预进制拓扑形式光滑”。


\subsec{第十五版补入的22.7相关正式勘误}
以下一组两项来自第四卷第四部分原分册第346页（上传PDF第345页），《勘误与增补（清单三）》A部；被订正页码指第四卷第一部分。

\medskip
\noindent\textbf{（$0_{\mathrm{IV}}$，22.7.3）。}
第212页倒数第21行，将 $\Ker(i)$ 改为 $\Ker(i_0)$。
第212页倒数第12行，将 $(\mathfrak p/\mathfrak p^2)\otimes_LL$ 改为 $(\mathfrak p/\mathfrak p^2)\otimes_AL$。


\subsec{第十六版补入的§23相关正式勘误}
以下两项来自第四卷第四部分原分册第346页（上传PDF第345页），《勘误与增补（清单三）》A部；被订正页码指第四卷第一部分。

\medskip
\noindent\textbf{（$0_{\mathrm{IV}}$，23.1.3）。}
第215页第4行，将 $x'^q\in K'$ 改为 $x'\in K'$。
第215页第9行，将 $x'^q\in A'$ 改为 $x'^q\in A$。
